% Bagian: propositional-logic
% Bab: syntax-and-semantics
% Subbagian: introduction

\documentclass[../../../include/open-logic-section]{subfiles}

\begin{document}

\olfileid[id]{pl}{syn}{int}

\olsection{Pendahuluan}

Logika proposisi membahas !!{formula}s yang dibangun dari
!!{propositional variable}s dengan menggunakan penghubung proposisional
$\lnot$, $\land$, $\lor$, $\lif$, dan $\liff$. Secara intuitif,
!!a{propositional variable}~$p$ mewakili sebuah kalimat atau proposisi
yang bernilai benar atau salah. Setiap kali ``nilai kebenaran''
!!{propositional variable} dalam !!a{formula} ditentukan, nilai kebenaran
setiap !!{formula}s yang dibentuk darinya dengan menggunakan penghubung
proposisional juga ditentukan. Kita mengatakan bahwa logika proposisi
bersifat \emph{fungsional-kebenaran}, karena semantiknya diberikan oleh
fungsi-fungsi atas nilai kebenaran. Secara khusus, dalam logika proposisi
kita tidak mempertimbangkan penentuan lebih lanjut mengenai kebenaran dan
kesalahan; misalnya, apakah sesuatu niscaya benar alih-alih sekadar benar
secara kontingen, apakah sesuatu diketahui benar, atau apakah sesuatu benar
sekarang alih-alih benar pada masa lalu atau masa mendatang. Kita hanya
mempertimbangkan dua nilai kebenaran, yaitu benar ($\True$) dan salah
($\False$), sehingga kita mengecualikan kemungkinan bahwa suatu pernyataan
tidak benar maupun tidak salah, atau hanya setengah benar. Kita juga hanya
memusatkan perhatian pada penghubung yang membuat nilai kebenaran
!!a{formula} yang dibangun dengannya sepenuhnya ditentukan oleh nilai
kebenaran bagian-bagiannya (dan bukan, misalnya, oleh maknanya). Secara
khusus, masih diperdebatkan apakah nilai kebenaran kalimat bersyarat dalam
bahasa Inggris bersifat fungsional-kebenaran dalam pengertian ini. Implikasi
material~$\lif$ memang demikian; logika-logika lain membahas kalimat
bersyarat yang tidak bersifat fungsional-kebenaran.

Untuk mengembangkan teori dan metateori logika proposisi berbasis fungsi
kebenaran, kita harus lebih dahulu mendefinisikan sintaks serta
semantik ungkapannya. Kita akan menguraikan salah satu cara
membangun !!{formula}s dari !!{propositional variable}s dengan menggunakan
penghubung-penghubung tersebut. Definisi alternatif juga dimungkinkan.
Sistem lain akan memilih simbol yang berbeda, memilih himpunan penghubung
primitif yang berbeda, dan menggunakan tanda kurung secara berbeda (atau
bahkan tidak menggunakannya sama sekali, seperti dalam apa yang disebut
notasi Polandia). Meskipun demikian, semua pendekatan memiliki kesamaan:
aturan pembentukan mendefinisikan himpunan !!{formula}s secara
\emph{induktif}. Jika dilakukan dengan tepat, setiap ungkapan pada dasarnya
hanya dapat dihasilkan melalui satu cara menurut aturan pembentukan.
Definisi induktif yang menghasilkan ungkapan yang \emph{terbaca secara
unik} memungkinkan kita memberikan makna kepada ungkapan tersebut dengan
metode yang sama---definisi induktif.

Pemberian makna kepada ungkapan merupakan ranah semantik. Konsep utama dalam
semantik logika proposisi ialah pemuasan dalam !!a{valuation}.
!!^a{valuation}~$\pAssign{v}$ memberikan nilai kebenaran $\True$, $\False$
kepada !!{propositional variable}s. Setiap !!{valuation} menentukan nilai
kebenaran $\pValue{v}(!A)$ bagi setiap !!{formula}~$!A$.
!!^a{formula} dipenuhi dalam !!a{valuation}~$\pAssign{v}$ jika dan hanya jika
$\pValue{v}(!A) = \True$---kita menuliskannya sebagai $\pSat{v}{!A}$.
Relasi ini juga dapat didefinisikan melalui induksi pada struktur~$!A$,
dengan menggunakan fungsi kebenaran bagi penghubung logis untuk
mendefinisikan, misalnya, pemuasan $!A \land !B$ berdasarkan dipenuhi atau
tidaknya $!A$ dan~$!B$.

Berdasarkan relasi pemuasan $\pSat{v}{!A}$ untuk kalimat, selanjutnya kita
dapat mendefinisikan gagasan semantis dasar berupa tautologi, perikutan, dan
keterpenuhan. !!^a{formula} merupakan tautologi, $\Entails !A$, jika setiap
!!{valuation} memenuhinya, yakni jika $\pValue{v}(!A) = \True$ untuk setiap
$\pAssign{v}$. Formula itu merupakan konsekuensi semantis dari suatu himpunan
!!{formula}s, $\Gamma \Entails !A$, jika setiap !!{valuation} yang memenuhi
semua !!{formula}s dalam~$\Gamma$ juga memenuhi~$!A$. Suatu himpunan
!!{formula}s dapat dipenuhi jika terdapat !!{valuation} yang memenuhi semua
!!{formula}s di dalamnya secara bersamaan. Karena !!{formula}s didefinisikan
secara induktif dan pemuasan selanjutnya didefinisikan melalui induksi pada
struktur !!{formula}s, kita dapat menggunakan induksi untuk membuktikan
sifat-sifat semantik kita serta menghubungkan gagasan-gagasan semantis yang
telah didefinisikan.


\end{document}
